\chapter{Covariance, coordinates and the geometry of the linear share}
\label{geometry:chapter}

There are two distinct ways to change a calculation involving several
traits. We can change the population while keeping the trait coordinates
fixed, or we can describe the same population in different coordinates.
The first can change the variance share. The second must leave it
unchanged if every quantity is transformed consistently. This chapter
proves both statements and explains the orientation comparisons used in
the visual examples.

Throughout, the population is Gaussian and the target is the specified
quadratic function $q$ from Chapter~\ref{gaussian:chapter}. A comparison
on a fixed surface also holds the population mean, gradient and Hessian
fixed unless a change is explicitly stated. The results are exact for
that quadratic function; applying them to a local expansion of a wider
surface retains the remainder qualifications in
Section~\ref{gaussian:taylor_derivation}.

\section{Proportional change in covariance}
\label{geometry:scaling}

Let $\mat G_0\succeq0$ be a reference covariance and define a family of
Gaussian populations with the same mean and covariance
$\mat G_s=s\mat G_0$, where $s>0$. This changes every variance and
covariance by the same factor. A standard deviation changes by $\sqrt{s}$;
the parameter $s$ is a variance multiplier, not a standard-deviation
multiplier. Put
\[
L_0=\vect b^{\mathsf T}\mat G_0\vect b,
\qquad
Q_0=\frac12\tr[(\mat H\mat G_0)^2].
\]

\begin{proposition}[Proportional-covariance scaling]
\label{geometry:scaling_proposition}
With the mean, $\vect b$ and $\mat H$ fixed, the two components under
$\mat G_s=s\mat G_0$ are
\begin{equation}
\Vlin(s)=sL_0,
\qquad \Vquad(s)=s^2Q_0.
\label{geometry:component_scaling}
\end{equation}
If $L_0+Q_0>0$, then
\begin{equation}
\Ag(s)=\frac{L_0}{L_0+sQ_0}.
\label{geometry:share_scaling}
\end{equation}
If both $L_0$ and $Q_0$ are positive, the share decreases strictly with
$s$, so proportional contraction increases it.
\end{proposition}
\begin{proof}
The linear expression contains one copy of the covariance:
$\vect b^{\mathsf T}(s\mat G_0)\vect b=sL_0$.
The quadratic expression contains two:
\[
\frac12\tr\{\mat H(s\mat G_0)\mat H(s\mat G_0)\}
 =s^2Q_0.
\]
Cancelling the common positive factor $s$ in the component ratio gives
Equation~\eqref{geometry:share_scaling}. Differentiating yields
\begin{equation}
\frac{d\Ag}{ds}
 =-\frac{L_0Q_0}{(L_0+sQ_0)^2},
\label{geometry:share_derivative}
\end{equation}
which is strictly negative precisely when both components are positive.
\end{proof}

The absolute quantities and their ratio therefore change at different
rates. For example,
\begin{align}
\frac{d\Vlin}{ds}&=L_0,
&\frac{d\Vquad}{ds}&=2sQ_0,\nonumber\\
1-\Ag(s)&=\frac{sQ_0}{L_0+sQ_0},
&\frac{d(1-\Ag)}{ds}&=\frac{L_0Q_0}{(L_0+sQ_0)^2}.
\label{geometry:omitted_fraction}
\end{align}
The omitted fraction does not scale as $s^2$, although the absolute
quadratic variance does. For $L_0>0$, its small-$s$ behaviour is
$1-\Ag(s)=sQ_0/L_0+O(s^2)$.

If $L_0,Q_0>0$, the mathematical limits are
\[
\lim_{s\downarrow0}\Ag(s)=1,
\qquad
\lim_{s\to\infty}\Ag(s)=0.
\]
At $s=0$ itself, however, both variances vanish and the share is undefined.
The first limit is not the share of a population with no variation.
If $L_0=0<Q_0$, the share remains zero for every positive $s$; if
$Q_0=0<L_0$, it remains one. Thus a narrow population at a stationary
point does not acquire a linear contribution merely by becoming narrower.
Large-$s$ limits describe the specified quadratic family, not an assurance
that a local expansion remains accurate over arbitrarily broad populations.

\paragraph{Biological interpretation.}
For a fixed local relationship and mean, narrowing the breeding-value
distribution can improve the fraction represented by a line while
reducing the amount of linearly described variation. There is no
contradiction: the numerator becomes smaller, but the curved contribution
becomes smaller faster. An actual comparison among populations must also
consider changes in the mean, conditional expected-fitness surface and
population-specific average effects. Holding these inputs fixed is how
this calculation isolates the frequency-distribution effect.

\subsection{Why an arbitrary contraction is different}
\label{geometry:anisotropic_counterexample}

The condition $\mat G_1\preceq\mat G_0$ says that variance has not
increased in any linear trait combination. It does not require all
directions to contract by the same fraction. Consider
\[
\vect b=\begin{pmatrix}1\\0\end{pmatrix},\qquad
\mat H=-\mat I_2,\qquad
\mat G_0=\mat I_2,\qquad
\mat G_1=\begin{pmatrix}0.1&0\\0&0.9\end{pmatrix}.
\]
The difference $\mat G_0-\mat G_1=\operatorname{diag}(0.9,0.1)$ is
positive definite. The initial components are
\[
V_{\mathrm{lin},0}=1,\qquad
V_{\mathrm{quad},0}=\tfrac12(1^2+1^2)=1,\qquad
\mathcal A_{g,0}=\tfrac12.
\]
After the contraction,
\[
V_{\mathrm{lin},1}=0.1,\qquad
V_{\mathrm{quad},1}=\tfrac12(0.1^2+0.9^2)=0.41,\qquad
\mathcal A_{g,1}=\frac{0.1}{0.51}\simeq0.1961.
\]
Both absolute components decreased, yet the share decreased too. Most of
the contraction was in the direction supplying the linear variation;
substantial variance remained in a direction supplying only curvature.
This example prevents extending Equation~\eqref{geometry:share_derivative}
to every covariance change described informally as a reduction in variation.

\section{An invertible change of trait coordinates}
\label{geometry:coordinate_change}

Let $\vect X'=\mat D\vect X$, where $\mat D$ is invertible. This is a
new description of the same individuals, not a new population. The
covariance becomes
\begin{equation}
\mat G'=\E[\vect X'\vect X'^{\mathsf T}]
       =\mat D\mat G\mat D^{\mathsf T}.
\end{equation}
The scalar target cannot change when its argument is relabelled. Since
$\vect X=\mat D^{-1}\vect X'$, its expression is
\begin{align}
q&=\ell_0+\vect b^{\mathsf T}\mat D^{-1}\vect X'
 +\frac12(\vect X')^{\mathsf T}
           \mat D^{-\mathsf T}\mat H\mat D^{-1}\vect X'.
\end{align}
Reading off its coefficients gives
\begin{equation}
\vect b'=\mat D^{-\mathsf T}\vect b,
\qquad
\mat H'=\mat D^{-\mathsf T}\mat H\mat D^{-1},
\label{geometry:transformation_rules}
\end{equation}
where $\mat D^{-\mathsf T}=(\mat D^{-1})^{\mathsf T}$.
Transforming a covariance while leaving the derivative coefficients in
their old units would describe a different calculation.

\begin{proposition}[Coordinate invariance]
\label{geometry:coordinate_invariance}
Under the simultaneous transformations above, the linear variance,
quadratic variance, quadratic mean contribution and defined share are
unchanged.
\end{proposition}
\begin{proof}
For the linear variance, adjacent inverse factors cancel:
\begin{align}
(\vect b')^{\mathsf T}\mat G'\vect b'
 &=\vect b^{\mathsf T}\mat D^{-1}
     (\mat D\mat G\mat D^{\mathsf T})
       \mat D^{-\mathsf T}\vect b\nonumber\\
 &=\vect b^{\mathsf T}\mat G\vect b.
\end{align}
For the quadratic term,
\begin{align}
\mat H'\mat G'
 &=\mat D^{-\mathsf T}\mat H\mat G\mat D^{\mathsf T},\nonumber\\
(\mat H'\mat G')^2
 &=\mat D^{-\mathsf T}\mat H\mat G
    (\mat D^{\mathsf T}\mat D^{-\mathsf T})
    \mat H\mat G\mat D^{\mathsf T}\nonumber\\
 &=\mat D^{-\mathsf T}(\mat H\mat G)^2\mat D^{\mathsf T}.
\end{align}
For compatible matrices, the trace identity
\[
\tr(\mat A\mat B)=\sum_{i,j}A_{ij}B_{ji}
                 =\tr(\mat B\mat A)
\]
allows a cyclic rotation of factors inside a trace. Applying it gives
\begin{align}
\tr[(\mat H'\mat G')^2]
 &=\tr\{(\mat H\mat G)^2
          \mat D^{\mathsf T}\mat D^{-\mathsf T}\}\nonumber\\
 &=\tr[(\mat H\mat G)^2].
\end{align}
The same argument with one copy of $\mat H'\mat G'$ proves invariance
of $\tr(\mat H\mat G)$. A cyclic rotation is not permission to rearrange
factors in an arbitrary order: matrices need not commute. The unchanged
components give an unchanged ratio whenever its denominator is positive.
\end{proof}

For a one-trait unit conversion $X'=10X$, the transformed quantities are
$G'=100G$, $b'=b/10$ and $H'=H/100$. Taking $G=1$, $b=2$ and $H=-0.5$
gives $G'=100$, $b'=0.2$ and $H'=-0.005$, and hence
\[
(0.2)^2(100)=4,\qquad
\tfrac12(-0.005)^2(100)^2=0.125.
\]
The numerical coefficients change, while the variances in log expected
output do not. The proof also applies to singular covariance matrices;
the coordinate transformation preserves their rank and maps the occupied
support to the same population expressed differently.

\section{Whitening and the curvature spectrum}
\label{geometry:whitening}

For $\mat G\succ0$, its spectral decomposition can be written
$\mat G=\mat U\mat\Lambda\mat U^{\mathsf T}$ with orthogonal
$\mat U$ and positive diagonal $\mat\Lambda$. Define the symmetric
square root and its inverse by
\[
\mat G^{1/2}=\mat U\mat\Lambda^{1/2}\mat U^{\mathsf T},
\qquad
\mat G^{-1/2}=\mat U\mat\Lambda^{-1/2}\mat U^{\mathsf T}.
\]
The rescaled variable $\vect Z=\mat G^{-1/2}\vect X$ has covariance
$\mat I$. For Gaussian $\vect X$, its density is proportional to
$\exp(-\sum_i z_i^2/2)$, a product of standard normal densities.
Thus the coordinates of $\vect Z$ are independent standard normals.
This covariance-removing transformation is called whitening.

Substitution $\vect X=\mat G^{1/2}\vect Z$ gives
\begin{equation}
q=\ell_0+\vect c^{\mathsf T}\vect Z
              +\frac12\vect Z^{\mathsf T}\mat S\vect Z,
\qquad
\vect c=\mat G^{1/2}\vect b,
\qquad
\mat S=\mat G^{1/2}\mat H\mat G^{1/2}.
\label{geometry:whitened_coefficients}
\end{equation}
The symmetric matrix $\mat S$ has real eigenvalues $\mu_i$. Write
$\mat S=\mat O\operatorname{diag}(\mu_i)\mat O^{\mathsf T}$ and
$\vect Y=\mat O^{\mathsf T}\vect Z$. Orthogonal rotation preserves a
standard normal density, so the $Y_i$ are again independent standard
normals. In these coordinates,
\begin{equation}
L=(\mat O^{\mathsf T}\vect c)^{\mathsf T}\vect Y,
\qquad
Q_c=\frac12\sum_i\mu_i(Y_i^2-1).
\end{equation}
Since $\Var(Y_i)=1$, $\Var(Y_i^2)=3-1=2$, and distinct squared
coordinates are independent, these forms give
\begin{equation}
\Vlin=\norm{\vect c}^2,
\qquad
\Vquad=\frac12\sum_i\mu_i^2.
\label{geometry:spectral_components}
\end{equation}
This also supplies a direct explanation of nonnegativity. Although
$\mat H\mat G$ need not be symmetric,
\begin{align}
\tr[(\mat H\mat G)^2]
 &=\tr[(\mat G^{1/2}\mat H\mat G^{1/2})^2]\nonumber\\
 &=\tr(\mat S^2)=\sum_i\mu_i^2\geq0.
\end{align}
For positive-definite $\mat G$,
$\mat H\mat G=\mat G^{-1/2}\mat S\mat G^{1/2}$, so it is similar
to $\mat S$. For singular $\mat G$, the whitening and similarity
statements are made on the occupied support; the trace identity itself
still holds in the full coordinates.

If $\norm{\vect c}$ and the eigenvalues $\mu_i$ are fixed, rotating
$\vect c$ relative to the eigenvectors of $\mat S$ cannot change the
share: neither scalar on the right of
Equation~\eqref{geometry:spectral_components} changes. The relative
orientation can affect selected-parent displacement, however, because
that calculation uses an inverse matrix acting on the gradient; see
Section~\ref{selection:eigenbasis_error}.

\paragraph{Biological interpretation.}
The number of named traits does not by itself determine how linear the
population description is. What matters here is the gradient and
curvature after accounting for the amount and association of variation in
the population. Adding a trait coordinate that never varies adds no
variance contribution. Relabelling the same varying coordinates cannot
change the share either.

\section{Physical orientation and the Rayleigh quotient}
\label{geometry:physical_orientation}

Now keep the trait axes fixed and change a biological input. For a unit
direction $\vect u$, the variance of the breeding-value combination
$\vect u^{\mathsf T}\vect X$ is
$\vect u^{\mathsf T}\mat G\vect u$. It is a directional evolvability
measure in the specified trait units \citep{HansenHoule2008}. Write a
gradient of fixed Euclidean length as $\vect b=a\vect u$, with $a>0$.
If $\mat G$ has eigenvectors $\vect v_i$ and eigenvalues $\lambda_i$,
expand $\vect u=\sum_i z_i\vect v_i$, where $\sum_i z_i^2=1$.
Then
\begin{equation}
\Vlin=a^2\vect u^{\mathsf T}\mat G\vect u
      =a^2\sum_i\lambda_i z_i^2.
\label{geometry:rayleigh_expansion}
\end{equation}
The weighted mean of the $\lambda_i$ lies between their minimum and
maximum. Hence
\begin{equation}
a^2\lambda_{\min}(\mat G)
 \leq\Vlin\leq a^2\lambda_{\max}(\mat G).
\label{geometry:rayleigh_bound}
\end{equation}
The upper equality holds when the gradient lies entirely in the
largest-eigenvalue eigenspace; if that eigenvalue is simple, this means
the corresponding eigenvector, with either sign. The lower equality is
analogous. This is the Rayleigh-quotient argument, including its equality
conditions. It fixes the gradient's ordinary length, not the length of
the whitened gradient.

Curvature's orientation relative to covariance can also matter. In a
two-trait coordinate system where $\mat G=\operatorname{diag}(g_1,g_2)$,
matrix multiplication gives
\begin{equation}
\Vquad=\frac12\{H_{11}^2g_1^2
                     +2H_{12}^2g_1g_2+H_{22}^2g_2^2\}.
\label{geometry:two_trait_curvature}
\end{equation}
Rotating a fixed set of curvature eigenvalues relative to these axes
changes the entries $H_{ij}$ and can change this variance. If instead
$\mat H=h\mat I$, then
$\Vquad=h^2\tr(\mat G^2)/2$, which depends only on the covariance
eigenvalues. A rotation of $\mat G$ with those eigenvalues fixed leaves
the quadratic variance unchanged.

This last case explains the two-trait rotation figure. Choose axes along
and perpendicular to the gradient, with gradient length $a$. Let the
population's larger-variance eigenvector make angle $\phi$ with the
gradient, and let its principal variances be $g_1$ and $g_2$. Resolving the
unit gradient direction along the two covariance eigenvectors gives
\begin{equation}
\Vlin=a^2(g_1\cos^2\phi+g_2\sin^2\phi),
\qquad
\Vquad=\frac{h^2}{2}(g_1^2+g_2^2)
\quad\text{when }\mat H=h\mat I.
\label{geometry:rotation_formula}
\end{equation}
At $\phi=90^\circ$, the linear variance is $a^2g_2$, not zero unless
the smaller principal variance is zero. A section along the long axis
can therefore have zero slope while the complete fitted plane still
varies in the perpendicular direction.

These physical comparisons do not contradict the fixed-spectrum statement
in Section~\ref{geometry:whitening}. Changing the relative position of
$\mat G$, $\vect b$ and $\mat H$ in fixed trait coordinates need not
preserve $\norm{\mat G^{1/2}\vect b}$ or the eigenvalues of
$\mat G^{1/2}\mat H\mat G^{1/2}$. Each comparison must state what stays
fixed before drawing an orientation conclusion.

\section{The regression triangle is in a space of functions}
\label{geometry:projection_triangle}

The geometry of a population regression is different from the physical
geometry of trait axes. For centred square-integrable functions $U$ and
$V$ of breeding value, define an inner product and length by
\[
\langle U,V\rangle=\E[UV],
\qquad \norm{U}_{L^2}=\sqrt{\E[U^2]}.
\]
Functions that agree almost surely are the same element for this
calculation. The population distribution enters the inner product:
changing its frequencies changes which differences matter to the fit.
For the Gaussian quadratic target,
\[
q-\E[q]=L+Q_c,
\qquad \langle L,Q_c\rangle=0.
\]
Expanding the square proves the right-triangle identity
\begin{equation}
\norm{q-\E[q]}_{L^2}^2
 =\norm{L}_{L^2}^2+\norm{Q_c}_{L^2}^2
 =\Vlin+\Vquad.
\label{geometry:pythagorean_projection}
\end{equation}
When $\Vlin>0$, define the angle $\psi$ between the centred target and
its linear projection through their normalised inner product. Then
\begin{align}
\cos\psi
 &=\frac{\langle L,L+Q_c\rangle}
         {\norm{L}_{L^2}\norm{L+Q_c}_{L^2}}\nonumber\\
 &=\frac{\Vlin}{\sqrt{\Vlin}\sqrt{\Vlin+\Vquad}},
\qquad \cos^2\psi=\Ag.
\label{geometry:angle_share}
\end{align}
The triangle's side lengths are standard deviations, not the component
variances themselves. It represents the relation among functions over a
population, not an angle between two trait vectors or a direction of
evolution. If $\Vlin=0<\Vquad$, the ordinary angle to a zero projection
is undefined; a plotted $90^\circ$ angle is only a limiting convention.
If both variances vanish, no such triangle is defined.

\paragraph{Biological interpretation.}
A large share states that the linear fit's mean squared residual is small
relative to the quadratic target's population variance. Neither a small function-space
angle nor a thin breeding-value ellipse specifies the descendants of the
population. Selection changes the weighting of the parental distribution,
and inheritance requires information beyond either geometric picture.

\section{Connection to executable checks}

\codepath{code/verify_mathematics.py} checks coordinate
transformations, proportional scaling and the explicit anisotropic
counterexample. The \codepath{orientation_and_absolute_scale} group in
\codepath{code/reproduce_result_checks.py} checks fixed
whitened norms separately from a physical-coordinate Rayleigh comparison.
\codepath{tutorials/geometry/code/geometry.py} integrates the
specified Gaussian populations to check the visual examples. These
finite computations illustrate the distinctions established in the
derivations; they do not substitute for the invariance or scaling proofs.
